\documentclass{article}

\usepackage{amsmath,amssymb}
\usepackage{xurl}
\usepackage[hidelinks]{hyperref}

% Shared commands for notes in docs/paper-gaps/.

\DeclareUnicodeCharacter{2080}{\ifmmode{}_{0}\else\textsubscript{0}\fi}
\DeclareUnicodeCharacter{2081}{\ifmmode{}_{1}\else\textsubscript{1}\fi}
\DeclareUnicodeCharacter{2082}{\ifmmode{}_{2}\else\textsubscript{2}\fi}
\DeclareUnicodeCharacter{2099}{\ifmmode{}_{n}\else\textsubscript{n}\fi}

\newcommand{\Lean}{\textsc{Lean}}
\ExplSyntaxOn
\NewDocumentCommand{\leanid}{m}{
  \ifmmode
    \textsf{\detokenize{#1}}
  \else
    \group_begin:
    \urlstyle{sf}
    \tl_set:Nn \l_tmpa_tl {#1}
    \tl_replace_all:Nnn \l_tmpa_tl {\_} {_}
    \exp_args:NV \path \l_tmpa_tl
    \group_end:
  \fi
}
\ExplSyntaxOff
\providecommand{\tr}{\operatorname{tr}}
\newcommand{\tnleanIssueBase}{https://github.com/LionSR/TNLean/issues/}
\newcommand{\tnleanPrBase}{https://github.com/LionSR/TNLean/pull/}
\newcommand{\tnleanDocBase}{https://sirui-lu.com/TNLean/docs/}
\newcommand{\ghissue}[1]{\href{\tnleanIssueBase#1}{GitHub issue \##1}}
\newcommand{\ghpr}[1]{\href{\tnleanPrBase#1}{PR \##1}}
\newcommand{\doclink}[1]{\href{\tnleanDocBase#1.html}{\path{#1}}}

% Dirac notation and the identity operator, as in blueprint/src/macros/common.tex.
% \providecommand keeps note-local definitions authoritative.
\usepackage{dsfont}
\providecommand{\ket}[1]{|#1\rangle}
\providecommand{\bra}[1]{\langle #1|}
\providecommand{\braket}[2]{\langle #1 | #2 \rangle}
\providecommand{\ketbra}[2]{|#1\rangle\!\langle #2|}
\providecommand{\Id}{\mathds{1}}
\providecommand{\one}{\mathds{1}}
\providecommand{\C}{\mathbb{C}}
\providecommand{\MN}[1]{M_{#1}(\C)}
\providecommand{\MD}{M_D(\C)}

% External referencing for paper sources.  The build helper writes sanitized
% auxiliary files under build/paper-aux/ and excludes duplicated source labels.
\usepackage{xr}

% Import a generated paper auxiliary file with a caller-supplied label prefix.
% The candidate paths support builds from docs/paper-gaps/, the repository root,
% and build/paper-aux/ itself.
\newcommand{\importpaperaux}[2]{%
  \IfFileExists{../../build/paper-aux/#2.aux}{%
    \externaldocument[#1]{../../build/paper-aux/#2}%
  }{%
    \IfFileExists{build/paper-aux/#2.aux}{%
      \externaldocument[#1]{build/paper-aux/#2}%
    }{%
      \IfFileExists{#2.aux}{%
        \externaldocument[#1]{#2}%
      }{}%
    }%
  }%
}

% General external reference macros
\newcommand{\paperref}[2]{\ref{#1-#2}}
\newcommand{\papereqref}[2]{(\ref{#1-#2})}

% CPSV16 source (1606.00608 / MPDO-22-12-17-2.tex) external document and shortcuts
\importpaperaux{cpsv16-}{1606.00608/MPDO-22-12-17-2-tnlean}
\newcommand{\cpsvref}[1]{\paperref{cpsv16}{#1}}
\newcommand{\cpsveqref}[1]{\papereqref{cpsv16}{#1}}

% Verdict marker: \gapnote{<kind>}{<status>}. Kind is the policy
% classification (clarification, local-correction, scope-restriction,
% unfaithful, false-source, open-gap); status is open, wip (elimination
% underway), resolved, or historical. Machine-read by the site index;
% prints a verdict line.
\newcommand{\gapnote}[2]{\par\noindent\textbf{Verdict:} #1 (#2).\par}


\title{Nilpotency Length One after Blocking: A Terminology Offset}
\author{}
\date{2026-09-03}

\begin{document}
\maketitle
\gapnote{clarification}{resolved}

\section{The two statements}

Let $A=(A^i)_i$ be an injective matrix product state tensor, let
$B=(B^i)_i$ be a tensor with the same physical alphabet and the same
positive-length closed-chain coefficients, and let $(V,W)$ be a reduction
from $B$ to $A$, so that $VW=\mathbf 1$ and
$VB^{i_1}\cdots B^{i_n}W=A^{i_1}\cdots A^{i_n}$ for every word.  Write
$N^i=B^i-WA^iV$ for the residual letters.  Definition~8 of
Ref.~\cite{Molnar2018SemiInjective} calls the least $N_0$ with
\begin{align}
    N^{i_1}\cdots N^{i_n}
        &=0
        \qquad (n\geq N_0)
    \label{eq:mgsc18_nilpotency_length_one_definition}
\end{align}
the nilpotency length of the selected reduction.\footnote{Source
    coordinates in \path{cornerproblem.tex} from
    \href{https://arxiv.org/abs/1706.07329v2}{arXiv:1706.07329v2}:
    lines~3147--3152.}
Thus $N_0=1$ means $N^i=0$ for every letter, that is, $B^i=WA^iV$.

Ref.~\cite{FrancoRubio2025MPUGauging} states the exterior fusion identity
\begin{align}
    B^{\mathbf p}WA^{\mathbf c}VB^{\mathbf q}
        &=B^{\mathbf p\mathbf c\mathbf q}
    \label{eq:mgsc18_nilpotency_length_one_exterior_identity}
\end{align}
for all nonempty central words $\mathbf c$ and all exterior words with
$|\mathbf p|,|\mathbf q|\geq\ell$, ``for some $\ell$ called the nilpotency
length'', and then says that ``$\ell$ can always be reduced to $1$ by
blocking the MPUs a finite number of times''.\footnote{\path{main.tex}
    line~1498 of \href{https://arxiv.org/abs/2502.20257}{arXiv:2502.20257};
    the fusion equations occupy lines~1409--1498.}
The footnote to Theorem~1 of Ref.~\cite{GarreRubioSchuch2024DomainWall}
makes the same assertion: the theorem ``assumes the nilpotency length of
the off-diagonal blocks of $B$ is $1$'', which ``can always be achieved by
blocking''.\footnote{\path{MPU-DW.tex} lines~358--364 of
    \href{https://arxiv.org/abs/2405.00439v2}{arXiv:2405.00439v2}.}

\section{The offset}

The corrected residual expansion
\cite{gap:mgsc18_residual_expansion_empty_word_error} gives
(\ref{eq:mgsc18_nilpotency_length_one_exterior_identity}) as soon as
\begin{align}
    N_0
        &\leq|\mathbf p|+1,
    &N_0
        &\leq|\mathbf q|+1,
    \label{eq:mgsc18_nilpotency_length_one_buffer_bounds}
\end{align}
because the adjacent central letter contributes one residual factor to
each exterior residual word.  Hence the exterior buffer length in the
sense of Ref.~\cite{FrancoRubio2025MPUGauging} is $\ell=N_0-1$ for
$N_0\geq1$, and ``$\ell=1$'' is the statement that one exterior site
suffices, that is, $N_0\leq2$.

Blocking $L$ original sites into one blocked site turns
(\ref{eq:mgsc18_nilpotency_length_one_buffer_bounds}) into the condition
$N_0\leq L+1$ for one blocked exterior site.  For a finite family of
reductions with nilpotency lengths $N_0(g,h)$, the length
\begin{align}
    L
        &=\max\bigl(1,\max_{(g,h)}(N_0(g,h)-1)\bigr)
    \label{eq:mgsc18_nilpotency_length_one_common_length}
\end{align}
is positive and satisfies $N_0(g,h)\leq L+1$ for every pair.  After this
common block, (\ref{eq:mgsc18_nilpotency_length_one_exterior_identity})
holds for every pair with $|\mathbf p|,|\mathbf q|\geq1$ blocked sites.
This is the content of the two quoted sentences that the formal
development uses.

Blocking does not, in general, make the blocked residual letters vanish,
and the blocked nilpotency length of
Definition~8 need not be $1$.  The two readings differ on nondegenerate
data.

\section{A counterexample to the literal reading}

Take the physical letters $i\in\{0,1\}$, the one-dimensional injective
tensor $A^0=A^1=1$, and the two-dimensional tensor
\begin{align}
    B^i
        &=\begin{pmatrix}
             1&x_i\\
             0&0
          \end{pmatrix},
    &x_0
        &=0,
    &x_1
        &=1.
    \label{eq:mgsc18_nilpotency_length_one_example_tensor}
\end{align}
Every word evaluates to $B^{i_1}\cdots B^{i_n}=B^{i_n}$, so all
positive-length closed-chain coefficients equal $1=A^{i_1}\cdots A^{i_n}$.
With $V=(1,0)$ and $W=(1,0)^{\mathsf T}$,
\begin{align}
    VW
        &=1,
    &VB^{i_1}\cdots B^{i_n}W
        &=1,
    &N^i
        &=\begin{pmatrix}
             0&x_i\\
             0&0
          \end{pmatrix},
    \label{eq:mgsc18_nilpotency_length_one_example_reduction}
\end{align}
so $(V,W)$ is a reduction, $N^1\neq0$, and $N^iN^j=0$.  Its nilpotency
length is $N_0=2$.

Block $L\geq1$ sites.  The blocked letter for the word
$\mathbf w=(i_1,\ldots,i_L)$ is $B^{\mathbf w}=B^{i_L}$, which depends on
the last original letter, while the blocked target letter is
$A^{\mathbf w}=1$.  For any reduction $(V',W')$ of the blocked tensors, the
blocked residual letter is $B^{i_L}-W'V'$, and the two values
$B^0-W'V'$ and $B^1-W'V'$ cannot both vanish.  Thus every reduction of
every blocking has nilpotency length $2$, and no blocking achieves $N_0=1$
in the sense of Definition~8.

One exterior site nevertheless suffices.  For $|\mathbf q|\geq1$,
\begin{align}
    B^{\mathbf p}WA^{\mathbf c}VB^{\mathbf q}
        &=B^{\mathbf p}\,WV\,B^{\mathbf q}
         =B^{i_{|\mathbf q|}}
         =B^{\mathbf p\mathbf c\mathbf q},
    \label{eq:mgsc18_nilpotency_length_one_example_exterior}
\end{align}
since $B^{\mathbf p}WV=B^0$ for every $\mathbf p$, and
$B^0B^{\mathbf q}=B^{\mathbf q}$ is the last letter of $\mathbf q$.  The
right buffer is necessary: for empty $\mathbf q$ and $\mathbf c=(1)$, the
left side is $B^{\mathbf p}WV=B^0\neq B^1=B^{\mathbf p\mathbf c}$.  This
agrees with (\ref{eq:mgsc18_nilpotency_length_one_buffer_bounds}) at
$N_0=2$.

\section{Convention}

The formal development reads ``the nilpotency length is reduced to $1$
by blocking'' as the common-block statement of
Section~2: one positive blocking length with $N_0(g,h)\leq L+1$ for every
pair, after which one blocked exterior site satisfies the exterior fusion
identity.  It does not require a blocked residual letter to vanish and
does not assert that the blocked nilpotency length equals $1$.\footnote{Formal
declarations:
\leanid{MPSTensor.IsReductionExteriorBufferLength},
\leanid{MPSTensor.IsReduction.blockTensor_commonBufferLength_one},
\leanid{MPOTensor.GroupFamily.ReductionFamily.block_hasExteriorBufferLength_one},
and
\leanid{MPOTensor.GroupFamily.IsRepresentation.exists_block_reductionFamily_one}.
The corresponding blueprint entries are
\path{thm:mps_reduction_common_blocking} in
\path{blueprint/src/chapter/ch02_mps.tex} and
\path{thm:mpug_common_blocking_one_site} in
\path{blueprint/src/chapter/ch29_mpu_gauging.tex}.}

\bibliographystyle{alpha}
\bibliography{references}

\end{document}
